\documentclass{beamer}

\usepackage[italian]{babel}
\usepackage[utf8]{inputenc}
\usepackage{makecell}
\usepackage{xcolor}

\usetheme{Madrid}

\title[Secure OTA su ESP32]
{Aggiornamenti OTA con Secure-Boot e Flash encryption sul dispositivo embedded ESP32}

\author[Irimie Fabio] {
  Irimie Fabio\texorpdfstring{\\}{, }
  Relatore: Prof. Fummi Franco
}

\institute[UniVr] {
  Università degli Studi di Verona\texorpdfstring{\\}{ - }
  Dipartimento di Informatica
}

\date
{Anno accademico 2023/2024}

\begin{document}

\frame{\titlepage}

\begin{frame}
\frametitle{Indice}
\tableofcontents
\end{frame}

\section{Obiettivo}
\begin{frame}
\frametitle{Obiettivo}
\begin{itemize}
  \item Proteggere il dispositivo da accessi non autorizzati alla memoria e dalla
    manomissione software tramite le seguenti funzionalità:

    \begin{itemize}
      \item \textbf{Flash Encryption:} Protegge i dati memorizzati nella memoria flash
        del dispositivo.
      \item \textbf{Secure Boot:} Impedisce l'esecuzione di firmware non autorizzato
      \item \textbf{Secure OTA (Over The Air):} Permette di aggiornare il firmware da
        remoto garantendo che il nuovo firmware sia autentico e non compromesso.
    \end{itemize}
\end{itemize}
\end{frame}

\section{Aggiornamenti OTA (Over The Air)}
\begin{frame}
\frametitle{Aggiornamenti OTA (Over The Air)} 

Gli aggiornamenti OTA permettono di aggiornare il firmware senza dover collegare il
dispositivo ad un computer.

\vspace{1em}
\noindent

Ci sono due modalità principali di aggiornamento e si distinguono in base alla
vulnerabilità del dispositivo.
\end{frame}

\subsection{Modalità di aggiornamento}
\begin{frame}
\frametitle{Modalità di aggiornamento}
\begin{itemize}
  \item \textbf{Modalità sicura:} Garantisce l'operabilità del dispositivo anche in
    caso di guasto durante l'aggiornamento. L'unico tipo di partizione che supporta
    questa modalità è:
    \begin{itemize}
      \item Application
    \end{itemize}
  \item \textbf{Modalità non sicura:} L'aggiornamento è vulnerabile a guasti e il
    dispositivo potrebbe non essere più operabile. Il tipo di partizioni che supportano
    questa modalità sono:
    \begin{itemize}
      \item Bootloader
      \item Partition Table
      \item Partizioni di tipo Data (come NVS, FAT, etc.)
    \end{itemize}
\end{itemize}
\end{frame}

\subsubsection{Aggiornamento in modalità sicura}
\begin{frame}
\frametitle{Aggiornamento in modalità sicura}
Vengono configurate le seguenti partizioni:
\begin{itemize}
  \item \texttt{ota\_0}
  \item \texttt{ota\_1}
  \item \texttt{ota\_data}: contiene informazioni sullo stato di boot
\end{itemize}
Durante l'aggiornamento, il nuovo firmware viene scritto nella partizione \texttt{ota}
non attiva. Se l'aggiornamento è completato con successo, la partizione \texttt{ota\_data}
viene aggiornata per indicare che la partizione \texttt{ota} appena scritta deve essere
utilizzata al boot successivo.

\vspace{1em}
\noindent
Se la partizione \texttt{ota\_data} non contiene
alcun dato il dispositivo esegue il boot dalla partizione \texttt{factory}.
\end{frame}

\subsubsection{Aggiornamento in modalità non sicura}
\begin{frame}
\frametitle{Aggiornamento in modalità non sicura}
Una partizione temporanea riceve
i dati della nuova immagine e, una volta completato il trasferimento, l'immagine viene
copiata nella partizione di destinazione.

\vspace{1em}
\noindent

Se l'operazione di copia viene interrotta potrebbero verificarsi problemi di boot.
\end{frame}

\subsection{App Rollback}
\begin{frame}
\frametitle{App Rollback}
L'App Rollback permette di tornare ad una versione precedente del firmware se la nuova
versione non funziona correttamente per garantire la funzionalità del dispositivo.
Il corretto funzionamento del dispositivo viene verificato tramite una funzione
eseguita dopo l'aggiornamento del firmware. 

\begin{itemize}
  \item Se l'app funziona correttamente, viene impostato lo stato dell'app a
    \texttt{ESP\_OTA\_IMG\_VALID}

  \item Se l'app non funziona correttamente, il dispositivo esegue il rollback alla versione
    precedente e imposta lo stato dell'app a \texttt{ESP\_OTA\_IMG\_INVALID}

  \item Se l’impostazione \texttt{CONFIG\_BOOTLOADER\_APP\_ROLLBACK\_ENABLE} è
    abilitata e viene effettuato un reset, allora viene effettuato un rollback senza
    chiamare nessuna funzione nell'applicazione.
\end{itemize}
\end{frame}


\section{eFuse}
\begin{frame}
  \frametitle{eFuse}
  Gli eFuse sono memorie non volatili che possono essere programmate una sola volta.
  Sono organizzati in 4 blocchi da 256 bit:
  \begin{itemize}
    \item \textbf{Blocco 0:} Riservato per valori di sistema
    \item \textbf{Blocco 1:} Chiave per Flash Encryption
    \item \textbf{Blocco 2:} Chiave per Secure Boot
    \item \textbf{Blocco 3:} Riservato per valori dell'utente
  \end{itemize}
\end{frame}

\section{Secure Boot}
\begin{frame}
\frametitle{Secure Boot}
Il Secure Boot verifica che ogni firmware eseguito sul dispositivo sia stato firmato
correttamente. Per la firma viene usata la crittografia assimmetrica:
\begin{itemize}
  \item \textbf{Chiave privata:} Mantenuta segreta e usata per firmare il firmware
  \item \textbf{Chiave pubblica:} Inserita negli eFuse del dispositivo per verificare
    l'autenticità del firmware
\end{itemize}
La firma è una struttura detta Signature Block (4Kb) che viene inserita alla fine
del firmware e viene calcolata su tutti i byte dell'immagine.
\end{frame}

\subsection{Processo di avvio con Secure Boot}
\begin{frame}
\frametitle{Processo di avvio con Secure Boot}
\begin{enumerate}
  \item Il first stage bootloader, presente in ROM, verifica la firma del second stage
    bootloader con la chiave presente negli eFuse.
  \item Se la firma è valida, il second stage bootloader viene eseguito,
    altrimenti il boot viene interrotto.
  \item Il second stage bootloader verifica la firma del firmware con la chiave
    presente negli eFuse.
  \item Se la firma è valida, il firmware viene eseguito,
    altrimenti si cerca un'altra immagine valida finchè non viene trovata o si
    esauriscono le immagini.
\end{enumerate}
\end{frame}

\section{Flash Encryption}
\begin{frame}
\frametitle{Flash Encryption}
La Flash Encryption permette di crittografare la memoria flash del dispositivo. 
Una volta attivata, tutte le seguenti partizioni vengono crittografate:
\begin{itemize}
  \item Second stage bootloader
  \item Partition table
  \item Partizione OTA Data
  \item Partizioni di tipo Application
\end{itemize}
Altre partizioni possono essere crittografate condizionalmente:
\begin{itemize}
  \item Partizioni con la flag \texttt{encrypted}
  \item Il hash del bootloader se il Secure Boot è abilitato
\end{itemize}
\end{frame}

\subsection{Modalità di Flash Encryption}
\begin{frame}
\frametitle{Modalità di Flash Encryption}
Esistono due modalità di funzionamento della Flash Encryption:
\begin{itemize}
  \item \textbf{Modalità sviluppo:} Permette di flashare nuove immagini non crittografate
    sul dispositivo e il bootloader le crittografa utilizzando la chiave memorizzata
    negli eFuse.

  \item \textbf{Modalità rilascio:} Non è possibile flashare nuove immagini non crittografate
    sul dispositivo, ma è necessario flashare immagini già crittografate.
\end{itemize}
\end{frame}

\subsection{eFuse rilevanti}
\begin{frame}
\frametitle{eFuse rilevanti}
I seguenti \texttt{eFuse} sono rilevanti per la Flash Encryption:
\begin{itemize}
  \item \texttt{flash\_encryption}: contiene la chiave di
    crittografia

  \item \texttt{DISABLE\_DL\_ENCRYPT} (1 bit): se abilitato disattiva la flash encryption
    mentre il dispositivo viene eseguito in modalità Firmware Download

  \item \texttt{DISABLE\_DL\_DECRYPT} (1 bit): se abilitato disattiva la flash decryption
    mentre il dispositivo viene eseguito in modalità UART Firmware Download

  \item \texttt{FLASH\_CRYPT\_CNT} (7 bit): indica se
    i contenuti della memoria flash sono stati crittografati.
    \begin{itemize}
      \item Se un numero dispari di bit sono impostati a 1, i contenuti
        della memoria flash sono crittografati e devono essere decrittografati durante
        la lettura.

      \item Se un numero pari di bit sono impostati a 1, i contenuti della
        memoria flash non sono crittografati.
    \end{itemize}
\end{itemize}
\end{frame}

\subsection{Processo di avvio con Flash Encryption}
\begin{frame}[allowframebreaks]
\frametitle{Processo di avvio con Flash Encryption}
Supponendo che i valori degli eFuse siano nello stato iniziale, il processo di avvio
con Flash Encryption è il seguente:
\begin{enumerate}
  \item Al primo avvio, la memoria flash non è crittografata e il first stage bootloader
    esegue il second stage bootloader.

  \item Il bootloader controlla il valore dell'eFuse
    \texttt{FLASH\_CRYPT\_CNT}. Siccome tutti i bit sono a 0 (stato di default),
    cioè un numero pari di bit impostati, il bootloader configura e attiva il
    \textit{flash encryption block}.

    \item Il bootloader utilizza la chiave memorizzata nell'eFuse
      \texttt{flash\_encryption}. Se la chiave non è presente ne genera una nuova e la
      scrive nell'eFuse.

  \item Il \textit{flash encryption block} crittografa tutti i dati nella memoria flash.
    L'intero processo avviene via hardware e la chiave non è accessibile
    dal software.

  \item Il bootloader imposta il primo bit disponibile di \\
    \texttt{FLASH\_CRYPT\_CNT} a 1 per indicare che i dati nella memoria flash sono
    crittografati. 

  \item In base alla modalità:
  \begin{itemize}
    \item In modalità \texttt{development}, il second stage bootloader abilita soltanto
      gli eFuse \texttt{DISABLE\_DL\_ENCRYPT} e \texttt{DISABLE\_DL\_CACHE} per
      permettere al bootloader in modalità Firmware Download di flashare binari crittografati.
      Inoltre i bit dell'eFuse \\ \texttt{FLASH\_CRYPT\_CNT} non sono protetti in scrittura.

    \item In modalità \texttt{release}, il second stage abilita i bit degli eFuse \\
      \texttt{DISABLE\_DL\_ENCRYPT}, \texttt{DISABLE\_DL\_DECRYPT} e \texttt{DISABLE\_DL\_CACHE}
      per impedire che il bootloader decrittografi i dati durante la modalità Firmware
      Download. Inoltre i bit dell'eFuse \texttt{FLASH\_CRYPT\_CNT} sono protetti in
      scrittura.
  \end{itemize}


  \item Il dispositivo viene riavviato per eseguire l'immagine crittografata. Il second
    stage bootloader chiama il \textit{flash decryption block} per decrittografare i dati
    della memoria flash e caricare i contenuti decrittografati in memoria.
\end{enumerate}
\end{frame}

\end{document}
